%% kapa: exam=1 round=28 year=2017 period=1 form=A subject=real_estate subject_ko=부동산학원론 seq=26 number=106
%% source: Q-Net 감정평가사 기출문제 (cst00309) artlSeq=5199603 file=extracted/1008/5199603/2017년 제28회 감정평가사 1차시험 문제지/1교시 A형.pdf
%% conversion: kordoc 4.24.51 -> tools/gen1.py (auto); review=proofread-sonnet; flags=none
\begin{question}{106}
감정평가유형에 관한 설명으로 옳지 않은 것은?
\begin{choices}
  \item 일괄평가란 2개 이상의 대상물건이 일체로 거래되거나 대상물건 상호간에 용도상 불가분의 관계가 있는 경우에는 일괄하여 평가하는 것을 말한다.
  \item 조건부평가란 일체로 이용되고 있는 물건의 일부만을 평가하는 것을 말한다.
  \item 구분평가란 1개의 대상물건이라도 가치를 달리하는 부분은 이를 구분하여 평가하는 것을 말한다.
  \item 현황평가란 대상물건의 상태, 구조, 이용방법 등을 있는 그대로 평가하는 것을 말한다.
  \item 참모평가란 대중평가가 아니라 고용주 혹은 고용기관을 위해 하는 평가를 말한다.
\end{choices}
\end{question}
\answer{2}
